\section{局部有限族}

记号约定:$X$是拓扑空间,$\left(A_{i}\right)_{i\in I}$是$X$的子集族.

\begin{definition}{局部有限族}{}
	若对每个$x\in X$,存在$x$的邻域$V$,使得$\left\{i\in I\middle|V\cap A_{i}\neq\emptyset\right\}$有限,则称$\left(A_{i}\right)_{i\in I}$是$X$的局部有限族.
\end{definition}

\begin{definition}{局部有限的子集}{}
	设$\mathcal{S}\subseteq\mathcal{P}\left(X\right)$.若$\mathcal{S}$构成的族是局部有限的,则称$\mathcal{S}$是$X$中的局部有限集.
\end{definition}

\begin{proposition}{局部有限族的子集族也局部有限,有限子集族局部有限}{}
	\begin{enumerate}
		\item 若$\left(A_{i}\right)_{i\in I}$局部有限,而对每个$i\in I$有$B_{i}\subseteq A_{i}$,则$\left(B_{i}\right)_{i\in I}$局部有限.
		\item $X$中的有限子集族都是局部有限族.
	\end{enumerate}
\end{proposition}

\begin{example}{局部有限但不有限的子集族}{}
	考虑$\R$中的开覆盖$\left(A_{n}\right)_{n=0}^{\infty}$:对任意$n\in\N$,有$A_{n}=
	\begin{cases}
		\left(-\infty,1\right),&n=0,\\
		\left(n-1,+\infty\right),&n\geq 1.
	\end{cases}$它显然是局部有限的,但并不是有限的.
\end{example}

\begin{proposition}{局部有限的闭集族的并是闭集}{}
	若$\left(F_{i}\right)_{i\in I}$是$X$的局部有限的闭集族,则$\bigcup\limits_{i\in I}F_{i}$是$X$的闭集.
\end{proposition}

\begin{example}{闭集族的并集不一定是闭集}{}
	拓扑空间$X$中,不是局部有限的闭集族的并不一定是闭集.例如,在$\Q$中,$\left(0,1\right)=\bigcup\limits_{n\geq 3}\left[\frac{1}{n},1-\frac{1}{n}\right]$,但是$\left(0,1\right)$不是闭集.
\end{example}